% !TEX root = main.tex
\subsection{Quantum Event Logic}
\label{sec:quantum}

In quantum logic, the events are the closed subspaces of a Hilbert
space~\cite{birkhoff-vonneumann1936}. A context is an orthogonal
basis, a maximal set of pairwise orthogonal rays. We take as events the
$18$ rays of Cabello's set in $\mathbb{R}^4$, spanned by vectors with
coordinates in $\{0, \pm 1\}$, and as contexts its $9$ orthogonal bases,
each ray lying in exactly two of them~\cite{cabello1996}\leanref{cabello-set}\leanref{basis-orth}.
We write $1\bar{1}00$ for the ray spanned by $(1, -1, 0, 0)$. A
\emph{state} is an element of the normalized local relaxation of
Definition~\ref{def:local}. For a unit vector $\psi$, the \emph{Born rule}
gives the ray spanned by $v$ the value $x_v = (\psi \cdot v)^2 / (v \cdot v)$. This is a state\leanref{born-state}\leanref{parseval}. The mixtures of these states are
the \emph{quantum states}.

The construction of Section~\ref{sec:calculi} gives the empty deterministic hull
here: no state has values in
$\{0,1\}$~\cite{kochen-specker1967,cabello1996}, so the normalized deterministic
hull is empty. We formalize this known
fact\leanref{no-ks}\leanref{no-certain-state}. We also prove that the
normalized local relaxation contains states that are not mixtures of
Born-rule assignments of these rays.

\begin{theorem}[local certainty in quantum logic\leanref{loc-quantum}]
\label{thm:loc-quantum}
On Cabello's set, the function $h$ with value $\tfrac12$ on the nine rays
$\bar{1}111$, $0001$, $0010$, $0100$, $1\bar{1}\bar{1}1$,
$1\bar{1}1\bar{1}$, $11\bar{1}1$, $111\bar{1}$, $1111$ and $0$ elsewhere is
a state and is not a mixture of Born-rule assignments.
\end{theorem}

\noindent
Every basis contains exactly two of the nine rays, so $h$ is a state. The
ray $0001$ lies in the plane of the orthogonal rays $1001$ and
$100\bar{1}$, so every Born-rule assignment satisfies
$x_{0001} \le x_{1001} + x_{100\bar{1}}$, the difference being
$\psi_1^2$. The function $h$ violates this law, with values $\tfrac12$,
$0$, and $0$. No basis contains $0001$ and either of the other two rays,
while each condition of the local relaxation involves a single basis.

In Theorem~\ref{thm:loc-quantum}, each ray $v$ is represented by the
projection onto $v$ in $\mathbb{R}^4$. A state $x$ has a \emph{quantum
model} if some Hilbert space has a projection $P_v$ for each ray $v$,
summing to the identity on each basis, and a positive operator
$\rho$ of trace $1$ with $x_v = \mathrm{tr}(\rho P_v)$ for every ray $v$.
Whether every state on Cabello's set has a quantum model was left
open~\cite{acin2015}. The state $h$ has none in finite dimension. It is
the only state that vanishes outside its nine rays, so it is an extreme point of the set of states. This step is not
formalized. An extreme state with a value outside
$\{0,1\}$ has no quantum model~\cite{sujan-kunjwal2025}. On all rays of $\mathbb{R}^d$
with $d \ge 3$, every state is a quantum state~\cite{gleason1957}.
